\documentclass[12pt]{article}

\usepackage[backref=page]{hyperref}
\usepackage{geometry}
\usepackage{amsmath}
\usepackage{amsfonts}
\usepackage{subcaption}
\usepackage{multirow}
\usepackage{parskip}
\usepackage{xcolor}
\usepackage{graphicx}
\usepackage{breakcites}
\usepackage{tikz}

\geometry{a4paper, margin=1in}
\graphicspath{{../..}}
\usetikzlibrary{arrows.meta}
\hbadness 10000
\vbadness 10000
\hfuzz 100pt
\vfuzz 100pt
\righthyphenmin 100
\emergencystretch 1em
\allowdisplaybreaks

\begin{document}

\begin{center}
    \Large
    \setlength{\baselineskip}{1.4em}
    \textbf{Algebraic Geometry For SLT: \\ Blow-Ups And Resolution Of Singularities} \\
    \vspace{12pt}
    \normalsize
    Jake Levi + ChatGPT $\cdot$ \today
\end{center}

\begin{abstract}
\noindent\begin{itemize}
    \item A singularity is often difficult to understand because several branches of a curve meet at the same point, or because a curve meets another curve in a complicated way.
    \item A \textbf{blow-up} is a geometric operation that replaces a point with a collection of directions, allowing us to separate these branches and examine how they meet.
    \item We can repeat this operation until the geometry becomes sufficiently simple.
    \item This process is called \textbf{resolution of singularities}.
    \item The reason this is useful for singular learning theory is that complicated polynomial functions can then be transformed into much simpler ones: locally, products of powers of individual coordinates.
    \item Such expressions are considerably easier to analyse, particularly when calculating integrals.
    \item The notes below build up the necessary intuition in the order in which you need it, using the cusp as a running example.
\end{itemize}
\end{abstract}

\setcounter{tocdepth}{1}
\tableofcontents

\newpage
\section{Vanishing locus (zero
set)}\label{vanishing-locus-zero-set}

The \textbf{vanishing locus} of a polynomial is simply the set of points
at which it evaluates to zero.

For example, consider a polynomial in two variables, \(f(a,b)=b^2-a^3\).
Its vanishing locus is

\[
V(f)=\{(a,b):b^2=a^3\}.
\]

Geometrically, this is a curve in the plane.

An important distinction: the vanishing locus tells us \emph{where} a
function is zero, but not \emph{how} it approaches zero. For example,
the univariate polynomials \(x\), \(x^2\) and \(x^3\) all have exactly the same vanishing locus,
namely the point \(x=0\), despite behaving very differently near that
point.

The behaviour of these polynomials is different in terms of \emph{how quickly they approach zero} near $x=0$: higher powers approach zero much faster. The important distinction is that these polynomials have the same zero set, but different multiplicities (or orders of vanishing). That's the information that the vanishing locus alone throws away.

This distinction becomes important when studying blow-ups.

The vanishing locus of a \emph{set} of polynomials is the set of points at which \emph{every} polynomial in the set is zero. It is the intersection of the vanishing loci of all polynomials in the set.

\newpage
\section{Ideal}\label{ideal}

An \textbf{ideal} is a collection of polynomials that is closed under
addition and multiplication by arbitrary polynomials.

Why introduce ideals? One reason is that they allow us to systematically
describe geometric objects (such as points, lines, curves, surfaces,
and intersections of these objects) using collections of polynomial
equations.

Formally, an \textbf{ideal} $I$ is a set of polynomials in which:
\begin{center}
    \begin{tabular}{p{0.3\linewidth} p{0.6\linewidth}}
        1.\; $f,g\in I \Rightarrow f+g\in I$
        & sums of polynomials in the ideal stay in the ideal, \\[0.7em]
        2.\; $(\forall h)\; f\in I \Rightarrow fh\in I$
        & products of a polynomial in the ideal and
        \emph{any other polynomial} stay in the ideal.
    \end{tabular}
\end{center}

For example, consider the set of all polynomials in $x$ that
vanish at both $x=3$ and $x=7$:
\[
I=\{f\in\mathbb{R}[x]:f(3)=f(7)=0\}.
\]
This is an ideal because adding two such polynomials still
gives zero at both points, and multiplying one by any other
polynomial also preserves these zeros. Its vanishing locus is $V(I)=\{3,7\}$.

\textbf{How do we specify an ideal?}
In the example above, we started with a geometric object
(the two points $\{3,7\}$) and defined an ideal by a
\emph{property}: the set of all polynomials that vanish at
those points.

But we can also start with
some polynomial equations and construct an ideal from them.
This is useful because we want to preserve information about
the equations themselves, not just their common zeros.

The \textbf{ideal generated by} a set of polynomials is the
smallest ideal containing them. In other words, it contains
the original polynomials and everything obtainable from them
by addition within the ideal and multiplication by arbitrary polynomials.
We denote this ideal by $(f_1,\ldots,f_n)$, where
$f_1,\ldots,f_n$ are the \textbf{generators}.

For example,
\[
I=(a,b)
\]
denotes the \textbf{ideal generated by} $a$ and $b$.
It contains $a$, $b$, $a+b$, $a^2$, $ab$, and all other
polynomial combinations of $a$ and $b$.

Its vanishing locus is
\[
V(I)=V(a)\cap V(b)=\{(0,0)\},
\]
because all polynomials in the ideal vanish exactly when
both generators vanish.

\textbf{Why retain the ideal rather than just its vanishing locus?}
Different ideals can have the same vanishing locus. For example,
\[
I=(a,b),\qquad J=(a^2,b^2)
\]
both vanish only at the origin, but they are different ideals:
$a\in I$ whereas $a\notin J$.

The ideal $J$ retains information about higher-order vanishing
that is lost when we consider only the vanishing locus.
This is precisely the information that \textbf{schemes},
introduced next, allow us to preserve.

\newpage
\section{Scheme}\label{scheme}

A \textbf{scheme} is a geometric object that retains not just the points
satisfying a collection of polynomial equations, but also the algebraic
information contained in those equations.

A useful way to think about it is:

\begin{itemize}
\item
  The vanishing locus tells us which points exist.
\item
  The scheme additionally tells us how the equations behave around those
  points.
\end{itemize}

For example, consider the following two schemes:

\[
X_1=(a=0), \qquad X_2=(a^2=0).
\]

Both have exactly the same underlying set of points: the line \(a=0\).

However, they are different schemes. The second retains the fact that
the equation vanishes to second order.

A particularly simple example is the \textbf{double point}:

\[
D=(a^2=0).
\]

Although its support consists of just one point, its functions behave as
though they have both a value and a first derivative at that point.
Algebraically, its ring of functions is

\[
k[\varepsilon]/(\varepsilon^2).
\]

Here \(\varepsilon\) is nonzero, but \(\varepsilon^2=0\).

You can think of this as a point with an infinitesimal direction
attached to it.

\textbf{Why this matters:} Schemes preserve information that would be
lost if we looked only at the underlying geometric set. In particular,
they preserve multiplicities and higher-order vanishing.

\newpage
\section{Reduced and non-reduced}\label{reduced-and-non-reduced}

A scheme is \textbf{reduced} if its equations contain no hidden
higher-order vanishing in the form of nonzero nilpotent elements.
Otherwise, it is \textbf{non-reduced}.

For example,

\[
(a=0)
\]

is reduced, whereas

\[
(a^2=0)
\]

is non-reduced.

Both describe the same underlying line, but the second describes it with
multiplicity two.

Similarly, the curve

\[
b^2=0
\]

is a line counted twice, rather than an ordinary line.

This distinction will reappear when we introduce divisors.

\newpage
\section{Singularities, smooth curves and
cusps}\label{singularities-smooth-curves-and-cusps}

A \textbf{smooth curve} is one that locally looks like an ordinary
straight line, possibly after changing coordinates.

A \textbf{singular point} is a point where this fails: the curve might
have a sharp point, intersect itself, or otherwise fail to look locally
like a smooth curve.

Two useful examples:

\textbf{Node (self-intersection):}

\[
ab=0.
\]

This describes the two coordinate axes, which intersect at the origin.
The origin is singular because two branches meet there.

\textbf{Cusp:}

\[
b^2=a^3.
\]

This describes a curve with a sharp point at the origin.

A parametrisation is

\[
a=t^2,\qquad b=t^3.
\]

Unlike the node, the cusp has just one branch. Nevertheless, that branch
has a singular point at the origin because it does not have a
well-defined smooth tangent direction there in the usual sense.

The cusp will be our main example when studying blow-ups.

\newpage
\section{Divisor and multiplicity}\label{divisor-and-multiplicity}

A \textbf{divisor} is a formal sum of curves, each assigned a
multiplicity.

For example,

\[
D=2\{a=0\}+3\{b=0\}
\]

describes the coordinate axes, with the first counted twice and the
second counted three times.

Why introduce this concept?

Because when we factorise a polynomial, we can interpret each factor as
a curve and its exponent as that curve's multiplicity.

For example,

\[
f(a,b)=a^2b^3
\]

corresponds to the divisor

\[
\operatorname{div}(f)
=2\operatorname{div}(a)+3\operatorname{div}(b).
\]

This is useful because it lets us separate two kinds of information:

\begin{itemize}
\item
  Which curves are present?
\item
  How many times does each curve occur?
\end{itemize}

Importantly, the divisor preserves multiplicities, unlike the vanishing
locus.

\textbf{The object we ultimately want to resolve is a divisor, not
merely a set of points.}

\newpage
\section{Simple normal crossings
(SNC)}\label{simple-normal-crossings-snc}

A divisor has \textbf{simple normal crossings} if, locally around every
point, its curves are smooth and intersect one another in the simplest
possible way.

In two dimensions, this means:

\begin{enumerate}
\def\labelenumi{\arabic{enumi}.}
\item
  Every individual curve is smooth.
\item
  When two curves intersect, their tangent directions are different.
\item
  No more than two curves meet at the same point.
\end{enumerate}

For example,

\[
D=\{a=0\}+\{b=0\}
\]

has simple normal crossings: two straight lines intersect at right
angles.

By contrast, the following situations do not have simple normal
crossings:

\begin{itemize}
\item
  Two curves touching tangentially rather than crossing.
\item
  Three curves meeting at one point.
\item
  A cusp.
\end{itemize}

The mathematical significance of SNC is that, in suitable local
coordinates, the divisor can be expressed as a monomial:

\[
D=\operatorname{div}(u_1^{m_1}u_2^{m_2}\cdots u_d^{m_d}).
\]

For example,

\[
f(u,v)=u^2v^3
\]

describes two coordinate axes, one with multiplicity two and the other
with multiplicity three.

This is precisely the kind of expression we want to obtain through
blow-ups. Monomials are particularly convenient because their integrals
can be separated into independent one-dimensional integrals.

\newpage
\section{Projective line}\label{projective-line}

The \textbf{projective line}, denoted \(\mathbb P^1\), is the collection
of all possible directions through the origin in a two-dimensional
plane.

A direction is represented by a pair of coordinates

\[
[x:y],
\]

where

\[
[x:y]=[cx:cy]
\]

for any nonzero constant \(c\).

In other words, \([1:2]\) and \([3:6]\) represent the same direction.

Why is this useful?

Ordinarily, a direction through the origin can be described by its
slope. But slope alone misses vertical lines. The projective line
includes \emph{all} directions, including vertical ones, in a single
mathematical object.

Over the real numbers, \(\mathbb P^1(\mathbb R)\) is topologically a
circle.

\textbf{Keep this intuition in mind:} the projective line is the
collection of all possible directions through a point. It is exactly
what will replace the origin in a blow-up.

\newpage
\section{Blow-up}\label{blow-up}

A \textbf{blow-up} is a geometric operation that replaces a point with
the projective line of directions through that point.

Consider the plane with coordinates \((a,b)\), and suppose we want to
blow up the origin.

The resulting space consists of pairs

\[
((a,b),[x:y])
\]

satisfying

\[
ay=bx.
\]

The equation simply says that the point \((a,b)\) lies on the line with
direction \([x:y]\).

There are two important cases.

\textbf{Away from the origin:} If \((a,b)\neq(0,0)\), there is exactly
one possible direction:

\[
[x:y]=[a:b].
\]

Consequently, nothing really changes away from the origin. The original
plane and the blown-up plane correspond exactly.

\textbf{At the origin:} If \((a,b)=(0,0)\), every direction satisfies
the equation.

Consequently, the origin is replaced by an entire projective line.

This new line is called the \textbf{exceptional divisor}.

\subsection{What does this achieve?}\label{what-does-this-achieve}

Imagine several straight lines crossing at the origin.

In the original plane, they all meet at the same point.

After blowing up the origin, each line meets the exceptional divisor at
a different point, corresponding to its direction.

The lines have therefore been separated!

This is the fundamental intuition behind blow-ups: replace a point where
several things are tangled together with the collection of directions
along which they approach that point.

\subsection{How do we actually calculate a
blow-up?}\label{how-do-we-actually-calculate-a-blow-up}

We work in two coordinate charts.

In the first chart, write

\[
a=u,\qquad b=uv.
\]

In the second chart, write

\[
a=rs,\qquad b=s.
\]

These are simply two ways of describing the same blown-up geometry. The
first chart covers all directions except the vertical one; the second
covers the remaining directions.

In the first chart, the exceptional divisor is

\[
u=0.
\]

Notice that \(u=0\) represents the entire collection of directions
through the original origin, rather than a single point in the original
plane.

These substitutions are the practical tools we use to calculate
blow-ups.

\newpage
\section{Exceptional divisor}\label{exceptional-divisor}

The \textbf{exceptional divisor} is the new curve introduced when we
blow up a point.

In the blow-up of the plane at the origin, it is the projective line

\[
E\simeq\mathbb P^1.
\]

Its points represent the different directions through the original
origin.

For example, the horizontal and vertical axes, which originally
intersected at the origin, now intersect the exceptional divisor at two
different points.

Importantly, when we perform subsequent blow-ups, we introduce
additional exceptional divisors. We must keep track of all of them,
because the final objective is to make the original curve and all the
exceptional divisors intersect with simple normal crossings.

\newpage
\section{Strict transform and total
transform}\label{strict-transform-and-total-transform}

Suppose we have a curve defined by

\[
f(a,b)=0
\]

and we blow up the origin.

After substituting

\[
a=u,\qquad b=uv,
\]

we obtain a new equation

\[
f(u,uv)=u^m\widetilde f(u,v),
\]

where \(m\) is the lowest total degree of any nonzero monomial in \(f\).

There are two things to distinguish.

\textbf{Strict transform:} This is the transformed version of the
original curve, excluding the exceptional divisor introduced by the
blow-up.

Its equation is simply

\[
\widetilde f(u,v)=0.
\]

Intuitively, it is the original curve, now viewed in the new geometry.

\textbf{Total transform:} This includes both the strict transform and
the exceptional divisor, with their respective multiplicities.

Its equation is the complete substituted equation

\[
f(u,uv)=u^m\widetilde f(u,v).
\]

For example, consider the node

\[
f(a,b)=ab.
\]

After substituting \(a=u,\ b=uv\),

\[
ab=u^2v.
\]

Therefore:

\begin{itemize}
\item
  The exceptional divisor is \(u=0\), with multiplicity two.
\item
  The strict transform is \(v=0\).
\end{itemize}

The total transform is thus

\[
\operatorname{div}(ab\circ\pi)=2E+\widetilde D.
\]

The original two branches have been separated: one is represented by the
strict transform, and the other is accounted for in the other coordinate
chart.

An especially useful fact is that the strict transform meets the
exceptional divisor precisely in the tangent directions of the original
curve.

\newpage
\section{Resolution of
singularities}\label{resolution-of-singularities}

\textbf{Resolution of singularities} is the process of repeatedly
applying blow-ups until the resulting divisor has simple normal
crossings.

The goal is not necessarily to make the original curve smooth in
isolation. Rather, we want the original curve, together with all the
exceptional divisors introduced along the way, to intersect in a simple,
predictable fashion.

\subsection{Example: resolving the
cusp}\label{example-resolving-the-cusp}

Consider

\[
f(a,b)=b^2-a^3.
\]

This is singular at the origin.

\textbf{First blow-up}

Substitute

\[
a=u,\qquad b=uv.
\]

Then

\[
f(u,uv)=u^2(v^2-u).
\]

The exceptional divisor is \(u=0\), and the strict transform is

\[
u=v^2.
\]

The strict transform is now smooth, but it is tangent to the exceptional
divisor. Therefore, we do not yet have simple normal crossings.

\textbf{Second blow-up}

Blow up the problematic intersection again, using

\[
u=rs,\qquad v=s.
\]

The strict transform becomes

\[
s^2-r=0.
\]

At this stage, three smooth curves meet at one point:

\begin{itemize}
\item
  The first exceptional divisor.
\item
  The second exceptional divisor.
\item
  The strict transform of the cusp.
\end{itemize}

Their tangent directions are different, but having three curves meet at
one point still violates the SNC condition.

\textbf{Third blow-up}

Blow up this triple intersection.

In a suitable chart, the resulting equations are

\[
E_3:\ u_3=0,
\]

\[
E_2:\ v_3=0,
\]

\[
\widetilde D:\ v_3-1=0.
\]

These curves now intersect with simple normal crossings. In the other
chart, the first exceptional divisor also intersects the third one
transversely.

We have therefore resolved the cusp!

The three blow-ups have progressively: 1. Made the strict transform
smooth, but left a tangency. 2. Removed the tangency, but created a
triple intersection. 3. Removed the triple intersection, achieving
simple normal crossings.

The final transformed divisor is

\[
\boxed{
\operatorname{div}(f\circ g)
=2E_1+3E_2+6E_3+\widetilde D.
}
\]

Notice that the multiplicities have been preserved throughout the
process.

\subsection{The general result}\label{the-general-result}

Hironaka's resolution theorem tells us that, for a polynomial, finitely
many suitable blow-ups suffice to obtain simple normal crossings.

In other words, however complicated the original polynomial's zero set
might be, we can transform it into something that locally behaves like a
monomial.

For singular learning theory, this is especially useful because the
transformed loss function and the Jacobian determinant of the
transformation can both be expressed locally as monomials times nonzero
factors.

This makes subsequent calculations, particularly asymptotic integrals,
much more tractable.

\newpage
\section{Glossary}

\begin{center}
    \begin{tabular}{p{0.2\linewidth} p{0.8\linewidth}}
        \raggedright\textbf{Vanishing locus} & The vanishing locus of a \emph{polynomial} is the set of zeros of the polynomial (where the polynomial ``vanishes''). The vanishing locus of a \emph{set} of polynomials is the \emph{intersection} of the vanishing loci of the polynomials in the set. \\[0.7em]
        \raggedright\textbf{Ideal} & An ideal is a set of polynomials which is closed under addition of polynomials in the ideal, and closed under multiplication of a polynomial in the ideal with any \emph{arbitrary} polynomial. Example: the set of polynomials which are zero at the points $\{3, 7\}$. An ideal not only defines a \emph{geometric} object (its vanishing locus, which may be a set of points or curve), but also, depending on which polynomials are or aren't in the ideal, implies \emph{algebraic} information about the multiplicity of the zeros of these polynomials. \\[0.7em]
        \raggedright\textbf{Scheme} & A scheme is a geometric object which is defined by a set of polynomials which are set equal to zero. A scheme has a \emph{support} and a \emph{ring of functions} (see below). \\[0.7em]
        \raggedright\textbf{Support (of a scheme)} & The support of a scheme is the vanishing locus of the set of polynomials. A scheme contains more information than its support, because the support does not contain information about the \emph{multiplicity} of zeros, which \emph{is} contained in the polynomials which define the scheme. \\[0.7em]
        \raggedright\textbf{Ring of functions (of a scheme)} & The ring of functions of a scheme is the set of \emph{all} polynomials in its variables (the variables of the polynomials defining the scheme), after the polynomials defining the scheme are ``applied'' (set equal to zero). In this ring of functions, many distinct polynomials become equivalent to each other, for example if they differ by a multiple of one of the equations defining the scheme (which is set equal to zero). \\[0.7em]
    \end{tabular}
\end{center}

\end{document}
